% !TEX root = main.tex

\section{Introduction}
\label{sec:introduction}

Over the last decade, Deep Learning has been successfully employed to tackle complex physical systems traditionally governed by non-linear differential equations. Central to this paradigm are Physics-Informed Neural Networks (PINNs)~\cite{raissi2019physics}, which directly embed the governing physical laws into the training loss via differential reconstruction residuals. By penalizing non-physical predictions, PINNs enable the discovery of continuous, mesh-free solutions even when observational data is sparse.

While standard PINNs have proven to be highly effective for purely differential problems, significant progress can be made extending these frameworks to systems with non-local memory. In this work, we focus on Volterra Integro-Differential Equations (VIDEs), a complex class of physical problems where the unknown state is governed not only by its immediate derivatives but also by its entire historical trajectory through an integral operator. This non-local memory dependence is a fundamental property governing critical real-world phenomena, such as epidemiology models, population dynamics, and the behavior of viscoelastic materials.

Existing deep learning solvers for VIDEs, such as the Residual Integration Solver Network (RISN)~\cite{moghaddam2025risn}, rely on Multi-Layer Perceptron (MLP) variants. However, standard MLPs operate under a strict \textit{pointwise} paradigm, processing spatial coordinates independently. Because Volterra operators are dominated by causal memory, these pointwise architectures are structurally ill-suited to capture non-local dependencies. Lacking internal relational mechanisms, they are forced to delegate the entire memory correlation to the external loss formulation.

Motivated by recent advances in incorporating attention mechanisms into PINNs and operator learning (e.g., Physics-Guided Transformers~\cite{zeraatkar2026pgt}, PINNsFormer~\cite{zhao2023pinnsformer}, AttnPINNs~\cite{song2026attnpinns}, and operator-level attention frameworks~\cite{kissas2022learning}), we investigate whether attention-based models can better address VID equations. Nevertheless, standard Transformer architectures rely on full Self-Attention encoders,
which suffer from severe computational complexity that dramatically restricts their applicability in continuous physical problems. To overcome this computational bottleneck, we propose \textbf{CrossPINN (C-PINN)}~\cite{rossi2026crosspinn}, an efficient architecture that bypasses dynamic sequence encoding in favor of a direct, learnable global latent basis interlinked via Sequential Cross-Attention.
While established complementary techniques—such as ReLoBRaLo loss balancing, Gauss-Legendre quadrature, and positional encodings—are directly adopted from the literature, our primary focus remains on isolating and evaluating the cross-attention mechanism itself as a solver for VIDEs.
The primary contributions of our implementation are the following:
\begin{itemize}
    \item We designed and implemented \textbf{CrossPINN (C-PINN)}, a domain-wise Cross-Attention architecture tailored to model non-local memory in Volterra operators via a direct learnable global latent basis.
    \item We documented the architectural evolution of the framework, detailing the iterative design choices required to adapt standard attention mechanisms to the strict constraints of physics-informed solvers.
    \item We performed a systematic benchmark evaluation comparing C-PINN variants against baseline MLP architectures across representative Volterra problems.
\end{itemize}

This report is structured as follows. Section~\ref{sec:background} reviews the VIDE mathematical formulation and related work. Section~\ref{sec:processing_pipeline} details the query-based processing pipeline and spatial coordinate streams. Section~\ref{sec:learning_framework} presents the C-PINN architecture and optimization strategies. Finally, Section~\ref{sec:results} discusses the benchmark results, followed by concluding remarks in Section~\ref{sec:conclusions}.